\section{Algorithmic Applications}\label{sec:optimization}

The geometric reformulation of magic square constraints as continuous covariance conditions suggests a natural optimization approach. By defining a smooth loss function $\mathcal{L}(M) = \sum_{k=1}^{6} c_k(M)^2$ based on the six covariances from Theorem~\ref{thm:main}, we transform the discrete combinatorial search into a continuous energy minimization problem.

\subsection{The Hybrid Gradient-Projection Algorithm}

Our approach operates in two phases:

\textbf{Phase 1: Continuous Descent.} Starting from a random real-valued $4 \times 4$ matrix, the algorithm follows the analytic gradient of the six-covariance loss, descending the energy landscape toward low-energy regions. This phase effectively filters the search space, guiding trajectories toward basins with ``magic-like'' structural properties.

\textbf{Phase 2: Integer Projection.} Once the loss falls below a specified tolerance, we project the continuous solution onto the discrete space of permutations. We employ the Hungarian algorithm to find the nearest valid permutation of $\{1, \ldots, 16\}$ to the continuous matrix. This projection step ``snaps'' the approximate solution to an exact discrete solution.

\subsection{Empirical Results for $n=4$}

\begin{table}[H]
\centering
\caption{Benchmark comparison for $4 \times 4$ magic square generation (1000 trials). All results are empirical observations on a single test system.}
\label{tab:benchmark}
\begin{tabular}{lrrr}
\toprule
\textbf{Method} & \textbf{Success Rate} & \textbf{Mean Time (ms)} & \textbf{Relative Speed} \\
\midrule
Gradient Descent (ours) & 99\% & 12.4 & $73\times$ \\
Backtracking & 100\% & 901.7 & $1\times$ \\
Simulated Annealing & 85\% & 45.3 & $20\times$ \\
\bottomrule
\end{tabular}
\end{table}

The gradient-based method achieves a 99\% success rate with mean convergence in approximately 280 iterations, representing a 73-fold empirical speedup over backtracking on our test system.

\begin{remark}[Scope and Limitations]
These results are purely empirical observations for $n=4$. We do not prove global convergence of the gradient descent procedure, polynomial-time complexity bounds, or that the optimization landscape possesses any convexity properties. The 1\% failure rate occurs when the continuous trajectory lands near a saddle point or local minimum that does not correspond to a magic square. For larger $n$, the landscape may become more rugged and success rates may decrease; extension to $n \geq 5$ remains untested.
\end{remark}

\subsection{Interpretation}

The empirical success of this hybrid method suggests that, at least for $n=4$, the basins of attraction around magic squares in the six-covariance landscape are sufficiently wide to capture random initializations with high probability. By relaxing discrete constraints to continuous ones, we can traverse the search space more efficiently---``sliding'' toward solutions rather than enumerating possibilities.

This result illustrates the potential of geometric relaxations for combinatorial problems, though it should be understood as a practical heuristic rather than a complexity-theoretic result.
